% ---
% title: Calling clients until the first sale
% tags: finance, geometric, expectation, variance
% ---

\begin{problem}
	\begin{context}
		A financial advisor calls potential clients one after another to sell
		shares of an investment fund. Each call either ends in a sale (the client
		buys) or not. We assume that every call ends in a sale with the same
		probability, independently of the other calls.
	\end{context}
	Each call is a sale with probability $0.2$, independently of the other
	calls. The advisor keeps calling until the first sale. Let $X$ be the number
	of the call on which the first sale happens.
	(a) What is the distribution of $X$, and what is its parameter?
	(b) Find $P(X=3)$.
	(c) Find $P(X>5)$, the probability that there is no sale in the first 5
	calls.
	(d) Find $\E[X]$ and $\SD(X)$.
	(e) The first 5 calls gave no sale. What is the probability that the next
	3 calls give no sale either? Compare with $P(X>3)$.
\end{problem}

\begin{slideproblem}
	Each call of an advisor is a sale with probability $0.2$, independently.
	$X$ = the number of the call with the first sale.
	(a) Distribution of $X$ and its parameter? \ (b) $P(X=3)$ \
	(c) $P(X>5)$ (no sale in the first 5 calls) \ (d) $\E[X]$, $\SD(X)$ \
	(e) The first 5 calls gave no sale. Probability that the next 3 give no
	sale either? Compare with $P(X>3)$.
\end{slideproblem}

\begin{slides}
	\small
	\exstatement
	\pause
	\begin{columns}[T]
		\begin{column}{.58\textwidth}
			\sol (a) Independent trials, same success probability, we wait for
			the \alert{first} success: $X$ is \alert{geometric} with $p=0.2$,
			\[ P(X=n)=0.8^{\,n-1}\cdot0.2,\qquad n=1,2,\dots \]
			\pause
			(b) $P(X=3)=0.8^2\cdot0.2=0.128$ \ (two failures, then a
			sale).\\[1mm]
			\pause
			(c) $P(X>5)=0.8^5\approx0.3277$ \ (the first 5 calls all fail).
		\end{column}
		\begin{column}{.4\textwidth}
			\centering
			\begin{tikzpicture}[x=.385cm,y=11cm]
			\draw[gray!40,very thin] (0.4,0.05)--(12.6,0.05);
			\node[left,font=\tiny] at (0.4,0.05) {$0.05$};
			\draw[gray!40,very thin] (0.4,0.1)--(12.6,0.1);
			\node[left,font=\tiny] at (0.4,0.1) {$0.1$};
			\draw[gray!40,very thin] (0.4,0.15)--(12.6,0.15);
			\node[left,font=\tiny] at (0.4,0.15) {$0.15$};
			\draw[gray!40,very thin] (0.4,0.2)--(12.6,0.2);
			\node[left,font=\tiny] at (0.4,0.2) {$0.2$};
			\draw[->] (0.4,0)--(12.9,0) node[right,font=\scriptsize] {$n$};
			\draw[->] (0.4,0)--(0.4,0.25) node[above,font=\scriptsize] {$P(X=n)$};
			\filldraw[fill=teal!22,draw=teal] (1-.38,0) rectangle (1+.38,0.2);
			\node[below,font=\tiny] at (1,0) {$1$};
			\filldraw[fill=teal!22,draw=teal] (2-.38,0) rectangle (2+.38,0.16);
			\node[below,font=\tiny] at (2,0) {$2$};
			\filldraw[fill=gold!55,draw=gold] (3-.38,0) rectangle (3+.38,0.128);
			\node[below,font=\tiny] at (3,0) {$3$};
			
			\filldraw[fill=teal!22,draw=teal] (4-.38,0) rectangle (4+.38,0.1024);
			\node[below,font=\tiny] at (4,0) {$4$};
			\filldraw[fill=teal!22,draw=teal] (5-.38,0) rectangle (5+.38,0.0819);
			\node[below,font=\tiny] at (5,0) {$5$};
			\filldraw[fill=navy!30,draw=navy] (6-.38,0) rectangle (6+.38,0.0655);
			\node[below,font=\tiny] at (6,0) {$6$};
			\filldraw[fill=navy!30,draw=navy] (7-.38,0) rectangle (7+.38,0.0524);
			\node[below,font=\tiny] at (7,0) {$7$};
			\filldraw[fill=navy!30,draw=navy] (8-.38,0) rectangle (8+.38,0.0419);
			\node[below,font=\tiny] at (8,0) {$8$};
			\filldraw[fill=navy!30,draw=navy] (9-.38,0) rectangle (9+.38,0.0336);
			\node[below,font=\tiny] at (9,0) {$9$};
			\filldraw[fill=navy!30,draw=navy] (10-.38,0) rectangle (10+.38,0.0268);
			\node[below,font=\tiny] at (10,0) {$10$};
			\filldraw[fill=navy!30,draw=navy] (11-.38,0) rectangle (11+.38,0.0215);
			\node[below,font=\tiny] at (11,0) {$11$};
			\filldraw[fill=navy!30,draw=navy] (12-.38,0) rectangle (12+.38,0.0172);
			\node[below,font=\tiny] at (12,0) {$12$};
			\node[font=\tiny,anchor=west] at (4.5,.2) {\textcolor{gold}{$\blacksquare$} (b) $X=3$: $0.128$};
				\node[font=\tiny,anchor=west] at (4.5,.165) {\textcolor{navy!60}{$\blacksquare$} (c) $X>5$: $0.328$};
			\end{tikzpicture}
		\end{column}
	\end{columns}
\end{slides}

\begin{slides}
	\small
	(d) Geometric with $p=0.2$:
	\[ \E[X]=\frac1p=\frac1{0.2}=5,\qquad
		\Var(X)=\frac{1-p}{p^2}=\frac{0.8}{0.04}=20,\qquad
		\SD(X)=\sqrt{20}\approx4.472. \]
	On average the first sale comes on the 5th call.
	\pause
	\medskip

	(e) No sale in calls 6, 7, 8 given no sale in calls 1--5:
	\[ P(X>8\mid X>5)=\frac{P(X>8)}{P(X>5)}
		=\frac{0.8^8}{0.8^5}=0.8^3=0.512. \]
	\pause
	And $P(X>3)=0.8^3=0.512$: the same!\\[1mm]
	\alert{Memoryless property:} after 5 failed calls the waiting starts
	afresh; past failures do not make a sale ``due''.
\end{slides}

\begin{solution}
	\idea The calls are independent trials with the same success probability
	$p=0.2$, and $X$ counts the trials up to and including the first success.
	This is exactly the setting of the geometric distribution, so every part is
	one of its formulas.

	\step{(a)} A ``success'' is a sale. The calls are independent and each one
	is a success with the same probability $0.2$. The number of the trial on
	which the first success happens is a geometric random variable, so $X$ is
	geometric with parameter $p=0.2$. Its pmf is
	\[ P(X=n)=(1-p)^{n-1}p=0.8^{\,n-1}\cdot0.2,\qquad n=1,2,3,\dots \]
	The reason: $X=n$ means that the first $n-1$ calls all fail (each with
	probability $0.8$) and call $n$ is a sale (probability $0.2$). The calls
	are independent, so these probabilities multiply.

	\step{(b)} $X=3$ means: no sale on call 1, no sale on call 2, sale on
	call 3. So
	\[ P(X=3)=0.8\cdot0.8\cdot0.2=0.8^2\cdot0.2=0.64\cdot0.2=0.128. \]

	\step{(c)} $X>5$ means that the first sale comes after call 5, that is,
	the first 5 calls are all failures. By independence,
	\[ P(X>5)=0.8^5=0.32768\approx0.3277. \]
	This is the formula $P(X>n)=(1-p)^n$. It is much shorter than adding up
	$1-\bigl(P(X=1)+\dots+P(X=5)\bigr)$, and gives the same result.

	\step{(d)} For a geometric random variable $\E[X]=1/p$ and
	$\Var(X)=(1-p)/p^2$. With $p=0.2$:
	\[ \E[X]=\frac{1}{0.2}=5,\qquad
		\Var(X)=\frac{0.8}{0.2^2}=\frac{0.8}{0.04}=20,\qquad
		\SD(X)=\sqrt{20}\approx4.472. \]
	On average the first sale comes on the 5th call. This makes sense: one call
	in five is a sale. The standard deviation is almost as large as the mean,
	so the waiting time varies a lot.

	\step{(e)} We know that the first 5 calls failed, that is, $X>5$. The next
	3 calls (calls 6, 7, 8) also fail exactly when $X>8$. We need the
	conditional probability $P(X>8\mid X>5)$. The event $X>8$ is contained in
	$X>5$, so the intersection of the two events is just $X>8$, and
	\[ P(X>8\mid X>5)=\frac{P(X>8)}{P(X>5)}=\frac{0.8^8}{0.8^5}=0.8^3=0.512. \]
	On the other hand, $P(X>3)=0.8^3=0.512$ as well: the probability that the
	first 3 calls give no sale. The two probabilities are equal.
	This is the \emph{memoryless property} of the geometric distribution: the
	past failures do not change the chances of the future calls, because the
	calls are independent. After 5 failed calls the advisor is in the same
	position as before the first call.
	\begin{answer}
		(a) geometric with $p=0.2$ \quad
		(b) $P(X=3)=0.8^2\cdot0.2=0.128$ \quad
		(c) $P(X>5)=0.8^5\approx0.3277$ \quad
		(d) $\E[X]=5$, $\Var(X)=20$, $\SD(X)=\sqrt{20}\approx4.472$ \quad
		(e) $P(X>8\mid X>5)=0.8^3=0.512=P(X>3)$
	\end{answer}
	\begin{remark}
		\textbf{Common mistake:} believing that after many failed calls a sale
		is ``due'', so that the next calls are more likely to succeed. With
		independent calls this is false: the memoryless property in (e) shows
		that the past failures change nothing. Another mistake is to use
		$P(X=n)=(1-p)^n p$: if the first sale is on call $n$, only the $n-1$
		calls before it fail.
	\end{remark}
\end{solution}
